\documentclass{article}

% --- PACKAGES ---
\usepackage[utf8]{inputenc}
\usepackage[T1]{fontenc}
\usepackage[hidelinks]{hyperref} % Removes the blue box/line from the link
\usepackage{geometry} % For setting margins
\usepackage{parskip} % Adds space between paragraphs and removes indent

% Customize margins for a cleaner letter look
\geometry{
    a4paper,
    top=2.5cm,
    bottom=2.5cm,
    left=2.5cm,
    right=2.5cm,
}

% Optional: Command to set the date manually (since the image shows Nov 30, 2025)
\newcommand{\LetterDate}{November 30, 2025}

\pagestyle{empty} % Removes page numbers

\begin{document}

% --- HEADER: RECIPIENT INFO (Left) and APPLICANT INFO (Right) ---
% Using two minipages to put them side-by-side
\noindent
\begin{minipage}[t]{0.5\textwidth} % Left half: Recipient Info
    ETH Zurich \\
    Admissions Office \\
    Rämistrasse 101 \\
    8092 Zürich \\
    Switzerland
\end{minipage}%
\hfill % Pushes the next minipage to the right
\begin{minipage}[t]{0.5\textwidth} % Right half: Applicant Info
    \begin{flushright}
        \textbf{Klevis Imeri} \\
        ANDRÁSSY ÚT 83-85 \\
        1062 BUDAPEST Hungary \\
        \href{mailto:klevisimeri11@gmail.com}{klevisimeri11@gmail.com} \\ % Mailto is better for email
        +36 70 538 4797
    \end{flushright}
\end{minipage}

\vspace{0.8cm} % Space between header and subject line

% --- SUBJECT LINE ---
\textbf{Motivation Letter for the MSc in Computer Science}

\vspace{0.5cm} % Space between subject and salutation

% --- SALUTATION ---
Dear Members of the Selection Committee,

\vspace{0.5cm}

% --- BODY TEXT (Final Spelling and Punctuation Corrections) ---
Since my first line of code in high school, where I was fed up doing
calculations by hand for the olympiads in physics, till
today, I have come to the conclusion that writing software is doable, but
writing reliable and secure software is hard. As a result, I have devoted
myself to researching ways to verify software without the need of human
labor.

One method to automate this process is model checking in formal methods.
That is the field I do my research on. Currently, I do this research at BME in the
FTSRG research group, where we apply and test our software in the Theta
framework using SV-COMP. I work on ways to pass the partial results from
different model checking algorithms like CEGAR, K-induction, and IMC. These
partial results happen when we decide on using some heuristics and AI models that
a specific algorithm is not able to solve this verification task. I have had
some good results on this, for which I also got first prize in the TDK (Student Scientific Conference) at BME. The reason I chose this field is that I love
mathematics. Especially exact mathematics with solvers and graphs. I must
mention, I just love a nice problem to think about. As a result, I also do a lot
of competitive programming with the Computing Programming Group at BME.
I also represented BME at ICPC CERC 2023.

When it comes to courses, I usually focus on security and networks, while I am
also part of the database group at BME as a lab instructor. I have a lot of
knowledge when it comes to internet protocols, especially in IPSec.

ETH Zurich has a lot of work done in this field, also especially in the Viper
Verification Infrastructure with its frontend tools. Formal methods had a
broader applicability than just verifying software, so there is also work in
applying these formal methods in security, especially verifying important
protocols with tools such as Tamarin Prover.

I dedicate my life to science and engineering. I spend nearly 12 hours every day
in this field. People say I am gifted. I only gift I have is being alive and
struggling towards something greater. I try to work the hardest to my humanly
capable limits. People usually judge my route for being inhuman routine. But this
constant struggle to reach the peak of research in an infinite large mountain, is what keeps me
going in life. It gives me the experience of the beauty of truth and the beauty
of creation.

Humans try to avoid change, because with change comes hardships to
adapt. True, I could also continue my master's at BME Budapest where I am familiar
with everything, and it would be quite easy. But I want to grow more, experience
more, no matter how hard it is. And ETH Zurich is the place for my next journey,
where it will allow me to do research in verifying software and also
expand the formal methods field in the networks and security.

\vspace{1cm} % Space between body and sign-off

% --- SIGN OFF: SIGNATURE (Left) and DATE (Right) ---
\noindent
\begin{minipage}[t]{0.5\textwidth} % Left half: Signature Block
    Sincerely,

    Klevis Imeri
\end{minipage}%
\hfill % Pushes the next minipage to the right
\begin{minipage}[t]{0.5\textwidth} % Right half: Date
    \raggedleft % Aligns the date to the right
    \LetterDate
\end{minipage}

\end{document}
